\chapter{Происхождение космологической связи из эйнштейновского отклика}
\label{ch:v10-cosmological-einstein-response}

\section{От энтропийного потока к тензору энергии}

Условный родитель предыдущего гейта связывал кривизну непосредственно с
производством энтропии. Эйнштейновское уравнение, однако, реагирует не на
безразмерную энтропию в единицу времени, а на плотность энергии--импульса.
Поэтому необходим явный промежуточный мост.

Пусть $\Theta$ --- энергия на единицу произведённой энтропии,
$\tau_{\rm res}$ --- время удержания этой энергии в одной ячейке, а
$v_{\rm cell}$ --- её физический четырёхмерный объём. Тогда условная плотность
энергии потока равна
\begin{equation}
 \rho_{\rm flow}
 =\frac{\Theta\sigma_\circlearrowright\tau_{\rm res}}{v_{\rm cell}}.
 \label{eq:v10-einstein-flow-density}
\end{equation}
Для одного времени проводимости
\begin{equation}
 \tau_{\rm res}=\frac1\kappa,
 \qquad
 \sigma_\circlearrowright=\kappa\log2,
\end{equation}
получаем
\begin{equation}
 \rho_{\rm flow}=\frac{\Theta\log2}{v_{\rm cell}}.
 \label{eq:v10-einstein-one-cycle-density}
\end{equation}

\section{Условная эйнштейновская калибровка}

Вакуумоподобный изотропный отклик имел бы космологическую кривизну
\begin{equation}
 \Lambda_E=\frac{8\pi G}{c^4}\rho_{\rm flow}
 =\frac{8\pi G\Theta\log2}{c^4v_{\rm cell}}.
 \label{eq:v10-einstein-curvature-response}
\end{equation}
Совмещение с динамической кривизной
$\Lambda_{\rm flow}=\kappa^2/(3c^2)$ даёт
\begin{equation}
 \boxed{\kappa^2=\frac{24\pi G\Theta\log2}{c^2v_{\rm cell}}}.
 \label{eq:v10-einstein-conductance-scale}
\end{equation}
В отличие от круговой формулы прошлого гейта это действительно могло бы
выбрать абсолютную проводимость, если бы $G$, $\Theta$ и $v_{\rm cell}$ были
получены независимо.

Три относительных условия можно собрать в положительный родитель. Для
нормировок
\begin{equation}
 u_\tau=\kappa\tau_{\rm res},
 \qquad u_E=\frac{\Lambda}{\Lambda_E},
 \qquad u_m=\frac{3c^2\Lambda}{\kappa^2}
\end{equation}
положим
\begin{equation}
 \mathcal P_E
 =\frac12(u_\tau-1)^2
 +\frac12(u_E-u_\tau)^2
 +\frac12(u_m-u_E)^2.
 \label{eq:v10-einstein-response-parent}
\end{equation}
Его единственный нуль равен $u_\tau=u_E=u_m=1$, а гессиан
\begin{equation}
 \operatorname{Hess}\mathcal P_E=
 \begin{pmatrix}
 2&-1&0\\
 -1&2&-1\\
 0&-1&1
 \end{pmatrix}
 \label{eq:v10-einstein-response-hessian}
\end{equation}
имеет ранг $3$, определитель $1$ и положительные ведущие миноры
$(2,3,1)$.

\section{Пакет недостающих якорей}

В логарифмических переменных $(\kappa,G,\Theta,v_{\rm cell})$ формула
\eqref{eq:v10-einstein-conductance-scale} даёт одну строку
\begin{equation}
 C_E=\begin{pmatrix}2&-1&-1&1\end{pmatrix}.
 \label{eq:v10-einstein-anchor-map}
\end{equation}
Её ранг равен $1$, а размерность ядра --- $3$. Независимая фиксация
$G$, $\Theta$ и $v_{\rm cell}$ повысила бы ранг полной карты до $4$.

Но текущая программа не располагает таким пакетом. Ньютоновская постоянная
не выведена из коэффициента эйнштейновского члена, КМС-температура наследует
невыведенную частоту часов, а абсолютный объём ячейки сохранил масштабную
моду. Кроме того, происхождение вакуумоподобного тензора энергии из
сквозного канала ещё не установлено.

\begin{theorem}[Условный масштаб из эйнштейновского отклика]
При независимых положительных $G$, $\Theta$ и $v_{\rm cell}$ совместный
родитель потока и кривизны выбирает
$\kappa^2=24\pi G\Theta\log2/(c^2v_{\rm cell})$. Его относительный гессиан
положительно определён и не содержит плоских направлений.
\end{theorem}

\begin{nogocriterion}[Энтропия не является энергией без моста]
Нельзя подставлять $\sigma_\circlearrowright$ непосредственно в правую часть
эйнштейновского уравнения. Для этого требуются энергия на единицу энтропии,
время удержания, абсолютный объём и типизированный тензор энергии. Пока их
происхождение не закрыто, формула масштаба остаётся условной.
\end{nogocriterion}

\begin{protocol}[Статус эйнштейновского моста]\begin{enumerate}
\item условный положительный родитель: \textbf{построен};
\item условная эйнштейновская калибровка: \textbf{$5/5$};
\item унаследованные источники: \textbf{$3/7$};
\item независимые $G$, $\Theta$, $v_{\rm cell}$ и тензор энергии:
\textbf{$0/4$};
\item абсолютная проводимость: \textbf{$0/1$};
\item следующий гейт: \textbf{аудит пакета якорей эйнштейновского отклика}.
\end{enumerate}\end{protocol}

Машинный сертификат сохранён в
\path{s2t_v10_cell_birth_four_volume_cosmological_constant_einstein_response_coupling_origin_gate_results.json}.